% !TEX root = ../notes.tex

\documentclass[../notes.tex]{subfiles}

\begin{document}

\section{May 11}
Today, we continue talking about flag manifolds.

\subsection{The Flag Manifold Is Projective}
Throughout, $G$ is a connected reductive group over $\CC$, and $\mc F$ is the flag manifold of Borel subgroups of $G$. Last time, we showed that $\mc F$ is compact. In fact, it is projective.
\begin{proposition} \label{prop:flag-proj}
	Fix a connected reductive group $G$ over $\CC$, and let $\mc F$ be the flag manifold. Then $\mc F$ is projective.
\end{proposition}
\begin{proof}
	We will embed $\mc F$ into $\PP V_\lambda$, where $V_\lambda$ is the finite-dimensional irreducible representation associated to a dominant integral weight $\lambda$. Notably, there is always a natural map $G\to\PP V_\lambda$ given by $g\mapsto[gv_\lambda]$, where $v_\lambda$ is the highest weight vector, so the difficulty will be to control the stabilizer of $v_\lambda$.
	
	To this end, choose a dominant integral weight $\lambda\in P^+$ which is regular in the sense that $\lambda(\alpha_i^\lor)\ge1$ for each simple coroot $\alpha_i^\lor$. Suppose that $L_\lambda$ lifts to an irreducible representation of $G$, and let $v_\lambda$ be the highest weight vector. Let $\mf b_+\subseteq\mf g$ be the standard Borel subalgebra. By definition of being highest weight, $v_\lambda$ is an eigenvector for $\mf b_+$ with eigenvalue $\lambda$. On the other hand, for each positive $\alpha$, we know $f_\alpha v_\lambda\ne0$: indeed,
	\[e_\alpha f_\alpha v_\lambda=(f_\alpha e_\alpha+h_\alpha)v_\lambda,\]
	which is $\alpha(\lambda)v_\lambda$, which is nonzero by our regularity. In fact, the vectors $f_\alpha v_\lambda$ have different weights, so they are linearly independent, so we conclude that
	\[\{X\in\mf g:Xv_\lambda\in\CC v_\lambda\}=\mf b_+.\]
	Thus, if $g\in G$ stabilizes $\CC v_\lambda$, then $g$ stabilizes $\mf b_+$, so \Cref{lem:borel-is-own-normalizer} implies that $g\in B_+$. It follows that the natural map $G\to\PP V_\lambda$ given by $g\mapsto[gv_\lambda]$ descends to an embedding $G/B_+\into\PP V_\lambda$.
\end{proof}
\begin{remark}
	The regularity condition is equivalent to $\lambda-\rho\in P_+$.
\end{remark}
\begin{remark}
	In fact, the map $G\to\PP V_\lambda$ in the argument is algebraic (because finite-dimensional representations are algebraic), so we have realized $\mc F$ as a complex projective variety.
\end{remark}
\begin{example}
	One can use this construction to recover the Pl\"ucker embedding of a Grassmannian into projective space.
\end{example}

\subsection{The Borel Fixed Point Theorem}
To continue, we need the following tool.
\begin{theorem}[Borel fixed point]
	Fix a solvable finite-dimensional Lie algebra $\mf a$, and let $V$ be a finite-dimensional representation of $\mf a$. If $X\subseteq\PP V$ is a closed $\mf a$-invariant subset, then $X$ admits a fixed point. Here, $\mc a$-invariance means that $e^ax\in X$ for any $a\in\mf a$ and $x\in X$.
\end{theorem}
\begin{proof}
	We proceed by induction on $\dim\mf a$. If $\mf a=0$, there is nothing to do. Otherwise, fix a proper nonzero ideal $\mf a'\subseteq\mf a$, which exists by the solvability of $\mf a$. Then note that $X^{\mf a'}$ is closed $\mf a$-invariant subset of $X$, so the induction gives it a fixed point. It now remains to find a point in $X^{\mf a'}$ which is fixed under the action of the quotient $\mf a/\mf a'$. Using solvability, we may assume that $\dim\mf a/\mf a'=1$, so we now have a one-dimensional Lie algebra acting on $X^{\mf a'}$.

	We have thus reduced to the case where $\dim\mf a=1$ acts on some closed invariant subspace $X\subseteq\PP V$. Here, $\mf a=\CC a$ for some $a\in\mf a$, so we are basically hunting for an eigenvector of $a$ in $X$. To this end, note that a change of basis of $V$ grants $a$ a Jordan normal form
	\[\bigoplus_iJ_{n_i}(\lambda_i).\]
	By scaling $a$ by a suitable complex number, we may assume that the real numbers $\Re\lambda_i$ are all distinct. We now appeal to the following linear algebraic lemma.
	\begin{lemma} \label{lem:limit-in-proj-space}
		Fix a complex linear operator $a\colon V\to V$ whose eigenvalues have distinct real parts. Then for each $x\in\PP V$, there is a limit
		\[\lim_{t\to\infty}e^{ta}x.\]
	\end{lemma}
	\begin{proof}
		Note that a change of basis of $V$ grants $a$ a Jordan normal form
		\[\bigoplus_iJ_{n_i}(\lambda_i).\]
		Now, by taking the direct sum decomposition, we may assume that $a$ is a single Jordan block $J_n(\lambda)$. Then
		\[\exp(tJ_n(\lambda))=\begin{bmatrix}
			e^{t\lambda} & te^{t\lambda} & \cdots & \frac{t^n}{n!}e^{t\lambda} \\
			& e^{t\lambda} & \ddots & \vdots \\
			&& \ddots & e^{t\lambda}\\
			&&& e^{t\lambda}
		\end{bmatrix}.\]
		We now see that the top-left coordinate gives the dominant term.
	\end{proof}
	To complete the proof, we choose any $x\in X$ and define $x_\infty\in\PP V$ to be the limit of $e^{ta}x$ as $t\to\infty$. Then $x_\infty$ is the desired fixed point; notably, it remains in $X$ because $X$ is $\mf a$-invariant and closed.
\end{proof}
\begin{remark}
	If $X=\PP V$, then this is Lie's theorem: we are hunting for an eigenvector for $\mf a$.
\end{remark}
\begin{remark}
	There is a similar statement for algebraic groups, whose proof is basically the same, except that no exponentiation is required.
\end{remark}

\subsection{Parabolic Subgroups}
Here is our definition.
\begin{definition}[parabolic]
	Fix a reductive Lie algebra $\mf g$. Then a Lie subalgebra $\mf p$ is \textit{parabolic} if and only if it contains a Borel subalgebra. If $G$ is a reductive Lie group with $\op{Lie}G=\mf g$, then the corresponding Lie subgroup $P$ is called a \textit{parabolic subgroup}.
\end{definition}
\begin{remark}
	Thus, a parabolic subgroup is one containing a Borel subgroup.
\end{remark}
\begin{remark} \label{rem:parabolic-subalgebra}
	Fix a polarization, which induces a Borel subgroup $\mf b_+\subseteq\mf g$. Then a parabolic subgroup $\mf p$ containing $\mf b_+$ admits a root decomposition
	\[\mf p=\mf h\oplus\bigoplus_{\alpha\in\Phi'}\mf g_\alpha,\]
	where $\Phi'\subseteq\Phi$ is some subset of roots. Notably, $\Phi'$ must contain the positive roots, but it may contain other ones. Additionally, as soon as $\Phi'$ contains some negative root $\beta$, because $\Phi'$ also contains all simple roots, we see that $\Phi'$ will contain every root greater than $\beta$.
\end{remark}
\begin{example}
	Fix $\mf g=\mf{gl}_n$, and let $\mf b_+$ be the Borel subalgebra of upper-triangular matrices. Then by choosing various negative roots as in \Cref{rem:parabolic-subalgebra}, we see that the parabolic subalgebras look like
	\[\mf b_++(\mf{gl}_{n_1}\oplus\cdots\oplus\mf{gl}_{n_k})\subseteq\mf{gl}_n,\]
	where $(n_1,\ldots,n_k)$ is an ordered partition of $n$. Thus, we see that $\mf p$ consists of ``block-upper-triangular matrices.'' 
\end{example}
Parabolic subalgebras, like all Lie algebras, admit a Levi decomposition. Namely, there is a canonical unipotent radical $U\subseteq P$, and one can even find splittings of the short exact sequence
\[1\to U\to P\to L\to1,\]
where $L\coloneqq P/U$. However, these splittings need not be unique in general.
\begin{remark}
	It turns out that all Levi subgroups inside $\mf p$ are conjugate by $P$. For example, if $P$ is just the Borel subgroup, then a Levi subgroup is exactly a maximal torus.
\end{remark}
\begin{example}
	The unipotent radical $\mf u$ of $\mf p$ consists of everything above the block-diagonal, and the Levi quotient $\mf p/\mf u$ admits a splitting to $\mf p$ given by exactly the block-diagonal matrices.
\end{example}
\begin{remark} \label{rem:parabolic-is-own-normalizer}
	As in \Cref{lem:borel-is-own-normalizer}, one finds that the normalizer of a parabolic subgroup is exactly the parabolic subgroup.
\end{remark}
We are now ready to define our partial flag manifolds.
\begin{definition}[partial flag manifold]
	Fix a complex reductive group $G$. For a parabolic subgroup $P\subseteq G$, the \textit{partial flag manifold} is $G/P$.
\end{definition}
\begin{remark}
	By \Cref{rem:parabolic-is-own-normalizer}, we see that $G/P$ parameterizes parabolic subgroups of $G$.
\end{remark}

\subsection{Maximal Solvable Subobjects}
As an application, let's run the following check.
\begin{proposition} \label{prop:borel-lie-subalgebra-maximal}
	Fix a complex reductive Lie algebra $\mf g$.
	\begin{listalph}
		\item Then $\mf b_+$ is a maximal solvable Lie subalgebra.
		\item Any solvable Lie subalgebra is contained in a Borel subalgebra.
	\end{listalph}
\end{proposition}
\begin{proof}
	Here we go.
	\begin{listalph}
		\item Any parabolic Lie subalgebra $\mf p$ containing $\mf b_+$ is given by adding in some negative roots by \Cref{rem:parabolic-subalgebra}. Thus, if $\mf p\supsetneq\mf b$, then $\mf p$ contains some $\mf{sl}_{2\alpha}$ for some root $\alpha$, so $\mf p$ fails to be solvable.
		\item Choose a solvable Lie subalgebra $\mf a\subseteq\mf g$. Then $\mf a$ acts on the flag manifold $\mc F$, which we know to be projective by \Cref{prop:flag-proj}. Thus, the action of $\mf a$ on $\mc F$ admits a fixed point $\mf b\in\mc F$, which means that $e^a\mf b=\mf b$ for each $a\in\mf a$, which means that $\mf a\subseteq\mf b$ because $\mf b$ is its own normalizer by \Cref{lem:borel-is-own-normalizer}.
		\qedhere
	\end{listalph}
\end{proof}
One can lift this to groups as follows.
\begin{proposition}
	Fix a complex reductive Lie group $G$.
	\begin{listalph}
		\item Then $B_+$ is a maximal solvable subgroup.
		\item Any solvable subgroup of $G$ is contained in a Borel subgroup.
	\end{listalph}
\end{proposition}
\begin{proof}
	Here we go.
	\begin{listalph}
		\item Fix a parabolic subgroup $P$ strictly containing $B_+$. Then we can find some $g\in P\setminus B_+$, meaning that $g$ does not normalize $\mf b_+$. Thus, $\op{Lie}P$ contains $\mf b_+$ and also contains the distinct Lie subalgebra $\op{Ad}_g\mf b_+$, so $\op{Lie}P$ is a parabolic Lie subalgebra strictly containing $\mf b_+$, so we are done by \Cref{prop:borel-lie-subalgebra-maximal}(a).
		\item Argue as in \Cref{prop:borel-lie-subalgebra-maximal}(b).
		\qedhere
	\end{listalph}
\end{proof}
\begin{example}
	For $\mf g=\mf{gl}_n$, we are finding that every solvable Lie subalgebra of $\mf{gl}_n$ can be upper-triangularized, which is Lie's theorem.
\end{example}

\subsection{Nilpotent and Unipotent Subobjects}
We can also say something about nilpotent subalgebras.
\begin{definition}[nilpotent subalgebra]
	A Lie subalgebra $\mf a\subseteq\mf g$ is a nilpotent subalgebra if and only if it acts by nilpotent operators on $\mf g$.
\end{definition}
\begin{remark}
	By Engel's theorem, nilpotent subalgebras are nilpotent Lie algebras, but the converse need not be true.
\end{remark}
\begin{corollary} \label{cor:nilp-subalg}
	Fix a reductive complex Lie algebra $\mf g$. Every nilpotent subalgebra of $\mf g$ is conjugate to a subalgebra of $\mf n_+$.
\end{corollary}
\begin{proof}
	Fix a nilpotent subalgebra $\mf a$. Then $\mf a$ can be conjugated into $\mf b_+$ by \Cref{prop:borel-lie-subalgebra-maximal}, so we may assume that $\mf a\subseteq\mf b_+$. However, $\mf a$ must be contained in $[\mf b_+,\mf b_+]$ to be nilpotent, so it follows that $\mf a\subseteq\mf n_+$.
\end{proof}
\begin{remark}
	It follows from \Cref{cor:nilp-subalg} that the maximal nilpotent subalgebras of $\mf g$ are conjugate to $\mf n_+$.
\end{remark}
Similarly, we can talk about unipotent subgroups.
\begin{definition}[unipotent]
	Fix a complex reductive Lie group $G$. Then a subgroup $A\subseteq G$ is a \textit{unipotent subgroup} if and only if $\op{Lie}A\subseteq\op{Lie}G$ is a nilpotent subalgebra. 
\end{definition}
\begin{example}
	The subgroup $N_+\coloneqq\exp(\mf n_+)$ is a unipotent subgroup.
\end{example}
\begin{corollary}
	Fix a complex reductive Lie group $G$. Then every unipotent subgroup of $G$ is conjugate to a subgroup of $N_+$.
\end{corollary}
\begin{proof}
	This follows from \Cref{cor:nilp-subalg}.
\end{proof}
\begin{remark}
	Once again, we see that maximal unipotent subgroups are conjugate to $N_+$.
\end{remark}
\begin{remark}
	Note that $\mf b_+$ contains a unique maximal nilpotent subalgebra, which is $\mf n_+$. Thus, there is a bijection between Borel subalgebras and maximal nilpotent subalgebras, so we see that the flag variety also classifies maximal nilpotent subalgebras. Namely, the normalizer of $\mf n_+$ is still $B_+$!
\end{remark}

\subsection{The Bruhat Decomposition}
Here is our result.
\begin{theorem}[Bruhat decomposition] \label{thm:bruhat}
	Fix a complex reductive group $G$. Let $B\subseteq G$ be a Borel subgroup containing the maximal torus $H$, and let $W$ be the Weyl group $N(H)/H$. Then
	\[G=\bigsqcup_{w\in W}BwB.\]
\end{theorem}
We will not prove this theorem in full, but we can show that these pieces are at least disjoint.
\begin{proposition}
	Fix a complex reductive group $G$. Let $B\subseteq G$ be a Borel subgroup containing the maximal torus $H$, and let $W$ be the Weyl group $N(H)/H$. Then the cosets $\{BwB\}_{w\in W}$ are disjoint.
\end{proposition}
\begin{proof}
	We may assume that $B$ is the standard Borel $B_+$. Fix $w_1,w_2\in N(H)$, and suppose that $Bw_1B=Bw_2B$. Then we would like to show that $w_1H=w_2H$. Well, we know that there are $b_1,b_2\in B$ for which
	\[b_1w_1=w_2b_2.\]
	Now, choose a regular dominant integral weight $\lambda$, and we compute the action of this element on some highest weight vector $v_\lambda$ of $L_\lambda$.
	\begin{itemize}
		\item On the right-hand side, we get $\lambda(b_2)w_2v_\lambda$, so this is an extremal vector of weight $w_2\lambda$.
		\item On the left-hand side, $w_1v_\lambda$ is an extremal vector of weight $w_1\lambda$, so $b_1w_1v_\lambda$ equals $w_1\lambda(b_1)w_1v_\lambda$ plus possibly some terms of other weights.
	\end{itemize}
	Because $\lambda(b_1)$ and $\lambda(b_2)$ are nonzero (they are just some scalars in $\CC^\times$), it follows that $w_1\lambda=w_2\lambda$. Thus, because $\lambda$ is regular, we conclude that $w_1=w_2$.
\end{proof}
Anyway, let's apply \Cref{thm:bruhat}.
\begin{corollary}
	Fix a complex reductive group $G$. Let $B\subseteq G$ be a Borel subgroup containing the maximal torus $H$, and let $W$ be the Weyl group $N(H)/H$. Then any pair of Borel subgroups of $G$ are conjugate to the pair $(B,w(B))$, where $w$ is in the Weyl group and is uniquely determined.
\end{corollary}
\begin{proof}
	Choose a pair $(B_1,B_2)$ of Borel subgroups. Certainly, we may conjugate $B_1$ to $B$, and we can even assume that $B_2$ looks like $\op{Ad}_gB$. By the Bruhat decomposition, we may expand $g$ as $b_1wb_2$ for some $b_1,b_2\in B$ and $w\in W$. Up to another conjugation, we thus see that $\left(B,\op{Ad}_gB\right)$ is conjugate to $(B,w(B))$.

	Lastly, for uniqueness, we note that $(B,w(B))$ and $(B,w'(B))$ being conjugate would imply that there is some $b\in B$ with $bw(B)b^{-1}=w'(B)$, but this eventually implies that $w$ and $w'$ live in the same Bruhat cell and thus $w=w'$.
\end{proof}
\begin{corollary}
	Fix a complex reductive group $G$. Let $B\subseteq G$ be a Borel subgroup containing the maximal torus $H$, and let $W$ be the Weyl group $N(H)/H$. Then
	\[G/B=\bigsqcup_{w\in W}BwB/B.\]
\end{corollary}
\begin{proof}
	Take the quotient of \Cref{thm:bruhat} by $B$.
\end{proof}
This is analogous to the Schubert cell decomposition for $\op{GL}_n$, so we can generalize our construction.
\begin{definition}[Schubert cells]
	Fix a complex reductive group $G$, and choose a Borel subgroup $B\subseteq G$. For some $w\in W$, the \textit{Schubert cell} $C_w$ of $w$ is the double coset
	\[BwB/B\subseteq G/B.\]
\end{definition}
\begin{remark}
	It turns out that the Schubert cells are locally closed.
\end{remark}
\begin{lemma} \label{lem:cell-decomp-flag}
	Fix a complex reductive group $G$. Let $B\subseteq G$ be a Borel subgroup containing the maximal torus $H$, and let $W$ be the Weyl group $N(H)/H$. For each $w\in W$, one has $C_w\cong\CC^{\ell(w)}$.
\end{lemma}
\begin{proof}
	By the Bruhat decomposition, $C_w=B/(B\cap w(B))$. Indeed, $B$ acts transitively on $C_w$, and the stabilizer of $wB\in BwB/B$ is exactly $B\cap w(B)$. Now, write $B=NH$, where $N$ is the unipotent radical. Then because $H\subseteq B\cap w(B)$, we conclude that
	\[C_w\cong\frac{N}{N\cap w(B)}=\frac N{N\cap w(N)}.\]
	We can now undo the exponential, and we see that this quotient is an affine space of dimension equal to the number of positive roots which go to negative roots.
\end{proof}
\begin{remark} \label{rem:closure-cell}
	One can check that $\overline C_w$ is highly singular.
\end{remark}
\begin{proposition}
	Fix a complex reductive group $G$, and let $\mc F$ be the flag manifold. Choose a coefficient ring $\Lambda$. Then the cohomology groups of $\mc F$ are free over $\Lambda$, and the Poincar\'e polynomial is
	\[\sum_{i\ge0}\op{rank}_\Lambda\mathrm H^i(\mc F;\Lambda)\cdot q^i=\sum_{w\in w}q^{2\ell(w)}.\]
\end{proposition}
\begin{proof}
	Use the cell decomposition given by \Cref{lem:cell-decomp-flag}; note that we can produce a genuine cell decomposition because the closure of a $G$-orbit is a union of smaller-dimensional orbits, so we can extend to a map on the closed cells. All differentials vanish in the cellular cochain complex because the cells are all even-dimensional.
\end{proof}
\begin{remark}
	It turns out that
	\[\sum_{w\in W}q^{\ell(w)}=\prod_{i=1}^r[m_i+1]_q,\]
	where $\{m_1,\ldots,m_r\}$ are the exponents of $\mf g$.
\end{remark}

\subsection{The Iwasawa Decomposition}
We close the course by stating the Iwasawa decomposition.
\begin{theorem}[Iwasawa decomposition]
	Fix a real form $G_\theta$ of a complex reductive group $G$. Let $\mf a\subseteq\mf p_\theta$ be a maximal abelian Lie subalgebra, and choose a generic $a\in\mf a$. Now, define $\mf n_{a^+}$ to be the sum of eigenspaces of $\op{ad}_a$ with positive eigenvalues, and let $N_{a^+}$ be the corresponding subgroup. Then the multiplication map
	\[K_c\times A\times N_{\mf a^+}\to G_\theta\]
	is a diffeomorphism.
\end{theorem}

\end{document}